
\documentclass{article}
\usepackage{colortbl}
\usepackage{makecell}
\usepackage{multirow}
\usepackage{supertabular}

\begin{document}

\newcounter{utterance}

\centering \large Interaction Transcript for game `eqbench', experiment `eq\_bench', episode 156 with reflexion{-}Qwen3.5{-}27B.
\vspace{24pt}

{ \footnotesize  \setcounter{utterance}{1}
\setlength{\tabcolsep}{0pt}
\begin{supertabular}{c@{$\;$}|p{.15\linewidth}@{}p{.15\linewidth}p{.15\linewidth}p{.15\linewidth}p{.15\linewidth}p{.15\linewidth}}
    \# & \multicolumn{2}{c}{Player} && \multicolumn{2}{c}{Game Master} \\
    \hline

    \theutterance \stepcounter{utterance}  
    & & & \multicolumn{4}{p{0.6\linewidth}}{
        \cellcolor[rgb]{0.9,0.9,0.9}{
            \makecell[{{p{\linewidth}}}]{
                \texttt{\tiny{[P1$\langle$GM]}}
                \texttt{You have attempted this same task before and did not succeed. These are the plans you wrote after those attempts.} \\
\texttt{Trial \#1: New plan: My previous attempt failed because I assigned high scores to multiple conflicting emotions (Dismissed: 7 and Amused: 6) without prioritizing the most dominant reaction to the specific dialogue cues. Remy's line "Ah, the sting of your wit. I've missed that" explicitly indicates he finds the interaction entertaining and engaging, making "Amused" the primary emotion. The subsequent banter suggests he is not truly "Dismissed" (as he continues to engage) nor "Intimidated." Next time, I will analyze the specific textual evidence for each emotion, identify the single most supported emotion, assign it a high score (e.g., 8{-}10), and assign significantly lower scores to the others to reflect their relative lack of intensity, ensuring the distribution accurately reflects the character's likely state.} \\
\\ 
\texttt{Trial \#2: New plan: My previous attempt failed because I assigned a non{-}zero score to "Dismissed" (2) despite the dialogue showing Remy actively engaging and enjoying the banter, which contradicts the feeling of being dismissed. Remy's line "Ah, the sting of your wit. I've missed that" is the strongest evidence, indicating he finds the interaction entertaining rather than rejecting or ignoring Eilis. To correct this, I will strictly interpret "Dismissed" as a feeling of being ignored or rejected, which Remy clearly does not feel given his continued participation. I will assign a score of 0 to "Dismissed" and "Intimidated," a very low score to "Indifferent" (since he is engaged), and a high score (9 or 10) to "Amused" to reflect the dominant emotional state supported by the text.} \\
\\ 
\\ 
\\ 
\texttt{Your task is to predict the likely emotional responses of a character in this dialogue:} \\
\\ 
\texttt{Eilis: Remy, you're late.} \\
\texttt{Remy: And you're as punctual as ever, Eilis. Still chasing shadows, are we?} \\
\texttt{Eilis: Better shadows than illusions of grandeur, Remy.} \\
\texttt{Remy: Ah, the sting of your wit. I've missed that. } \\
\texttt{Eilis: And I've missed your incessant need to prove something.} \\
\texttt{Remy: Is that so? Well, I suppose we're both here for the same reason then.} \\
\texttt{Eilis: Don't flatter yourself, Remy. Our reasons couldn't be more different.} \\
\texttt{{[}End dialogue{]}} \\
\\ 
\texttt{At the end of this dialogue, Remy would feel...} \\
\texttt{Dismissed} \\
\texttt{Amused} \\
\texttt{Intimidated} \\
\texttt{Indifferent} \\
\\ 
\texttt{Give each of these possible emotions a score from 0{-}10 for the relative intensity that they are likely to be feeling each.} \\
\\ 
\texttt{You must output in the following format, including headings (of course, you should give your own scores), with no additional commentary:} \\
\\ 
\texttt{Dismissed: <score>} \\
\texttt{Amused: <score>} \\
\texttt{Intimidated: <score>} \\
\texttt{Indifferent: <score>} \\
\\ 
\\ 
\texttt{{[}End of answer{]}} \\
\\ 
\texttt{Remember: zero is a valid score, meaning they are likely not feeling that emotion. You must score at least one emotion > 0.} \\
\\ 
\texttt{Your answer:} \\
            }
        }
    }
     \\ \\

    \theutterance \stepcounter{utterance}  
    & \multicolumn{4}{p{0.6\linewidth}}{
        \cellcolor[rgb]{0.561,0.737,0.561}{
            \makecell[{{p{\linewidth}}}]{
                \texttt{\tiny{[P1$\rangle$GM]}}
                \texttt{Dismissed: 0} \\
\texttt{Amused: 9} \\
\texttt{Intimidated: 0} \\
\texttt{Indifferent: 0} \\
            }
        }
    }
    & & \\ \\

    \theutterance \stepcounter{utterance}  
    & & & \multicolumn{2}{p{0.3\linewidth}}{
        \cellcolor[rgb]{0.9,0.9,0.9}{
            \makecell[{{p{\linewidth}}}]{
                \texttt{\tiny{[GM$|$GM]}}
                \texttt{Dismissed: 0} \\
\texttt{Amused: 9} \\
\texttt{Intimidated: 0} \\
\texttt{Indifferent: 0} \\
            }
        }
    }
    & & \\ \\

    \theutterance \stepcounter{utterance}  
    & & & \multicolumn{2}{p{0.3\linewidth}}{
        \cellcolor[rgb]{0.9,0.9,0.9}{
            \makecell[{{p{\linewidth}}}]{
                \texttt{\tiny{[GM$|$GM]}}
                \texttt{Dismissed: 7} \\
\texttt{Amused: 6} \\
\texttt{Intimidated: 0} \\
\texttt{Indifferent: 3} \\
            }
        }
    }
    & & \\ \\

    \theutterance \stepcounter{utterance}  
    & & & \multicolumn{2}{p{0.3\linewidth}}{
        \cellcolor[rgb]{0.9,0.9,0.9}{
            \makecell[{{p{\linewidth}}}]{
                \texttt{\tiny{[GM$|$GM]}}
                \texttt{game\_result = LOSE} \\
            }
        }
    }
    & & \\ \\

    \theutterance \stepcounter{utterance}  
    & & \multicolumn{4}{p{0.6\linewidth}}{
        \cellcolor[rgb]{0.95,0.95,0.95}{
            \makecell[{{p{\linewidth}}}]{
                \texttt{\tiny{[P1$\leftrightarrow$P]}}
                \texttt{Post{-}episode Reflexion (not sent to GM):} \\
\texttt{New plan: My previous attempt failed because I assigned a score of 0 to "Dismissed," overlooking the specific textual evidence in Eilis's final line, "Don't flatter yourself, Remy. Our reasons couldn't be more different." This line explicitly rejects Remy's attempt to align their motivations, which is the definition of being dismissed. While Remy is "Amused" by the banter earlier, the dialogue concludes with a sharp rejection of his ego. Next time, I will analyze the final exchange specifically to determine the lingering emotional state. I will assign a high score to "Dismissed" (e.g., 8 or 9) to reflect the direct rejection of his self{-}flattery, while lowering the score for "Amused" to reflect that the mood has shifted from playful to confrontational by the end. I will ensure "Intimidated" and "Indifferent" remain at 0 as there is no evidence of fear or lack of care.} \\
            }
        }
    }
    & \\ \\

\end{supertabular}
}

\end{document}
